\documentclass{article}
\usepackage{amsmath,amssymb,amsthm}
\newtheorem{theorem}{Theorem}
\begin{document}

\begin{theorem}
On an interval where the eight operators
\[
srx,\ sxp,\ cxp,\ crx,\ urx,\ sawup,\ uxp,\ sawdown
\]
are defined and nonzero (for example \(0<x<\pi/2\)), set
\[
f_1=srx,\quad f_2=sxp,\quad f_3=cxp,\quad f_4=crx,\quad
f_5=urx,\quad f_6=sawup,\quad f_7=uxp,\quad f_8=sawdown.
\]
Define the \(9\times 9\) skew-symmetric matrices
\[
Y_i=f_i\,(E_{i,i+1}-E_{i+1,i})\qquad (i=1,\dots,8).
\]
Then the Lie algebra \(\mathfrak{g}=\langle Y_1,\dots,Y_8\rangle\) generated by
these matrices is \(\mathfrak{so}(9)\). Consequently,
\[
\dim\mathfrak{g}=36,\qquad \operatorname{rank}\mathfrak{g}=4.
\]
\end{theorem}

\begin{proof}
Fix \(x\) in the interval, so each \(f_i\) is a nonzero real number. Since
\(f_i\) is a unit, we may divide by it and obtain
\[
A_i \;:=\; E_{i,i+1}-E_{i+1,i} \;=\; f_i^{-1} Y_i \;\in\; \mathfrak{g}
\qquad (i=1,\dots,8).
\]

We show by induction on the distance \(k\) that every elementary skew matrix
\[
A_{i,i+k} \;:=\; E_{i,i+k}-E_{i+k,i}
\]
lies in \(\mathfrak{g}\), for \(1\le k\le 8\) and \(1\le i\le 9-k\).
The case \(k=1\) is \(A_{i,i+1}=A_i\in\mathfrak{g}\). Suppose
\(A_{i,i+k}\in\mathfrak{g}\) for some \(1\le k\le 7\). Bracketing with
\(A_{i+k}=E_{i+k,i+k+1}-E_{i+k+1,i+k}\) gives
\[
[A_{i,i+k},\,A_{i+k}]
= [E_{i,i+k}-E_{i+k,i},\,E_{i+k,i+k+1}-E_{i+k+1,i+k}]
= E_{i,i+k+1}-E_{i+k+1,i}
= A_{i,i+k+1},
\]
using \(E_{a,b}E_{c,d}=\delta_{b,c}E_{a,d}\): of the four products in each
order, only \(E_{i,i+k}E_{i+k,i+k+1}=E_{i,i+k+1}\) and
\(E_{i+k+1,i+k}E_{i+k,i}=E_{i+k+1,i}\) survive, and they enter with opposite
signs. Hence \(A_{i,i+k+1}\in\mathfrak{g}\), completing the induction.

Therefore \(\mathfrak{g}\) contains \(A_{i,j}=E_{i,j}-E_{j,i}\) for every
\(1\le i<j\le 9\). These \(\binom{9}{2}=36\) matrices are linearly independent
and span the space of real skew-symmetric \(9\times 9\) matrices, which is
\(\mathfrak{so}(9)\). Every \(Y_i\) is skew-symmetric, so
\(\mathfrak{g}\subseteq\mathfrak{so}(9)\), and the spanning gives
\[
\mathfrak{g}=\mathfrak{so}(9).
\]

Finally,
\[
\dim\mathfrak{so}(9)=\frac{9\cdot 8}{2}=36,\qquad
\operatorname{rank}\mathfrak{so}(9)=\left\lfloor\frac{9}{2}\right\rfloor=4,
\]
and the root system is of type \(B_4\).
\end{proof}

\end{document}
